\subsection{One resolved obstacle suffices for exact recovery}
\label{sec:single-resolved-extension}
The uniform width restriction can be removed from all but one
obstacle.  This extension concerns exact recovery; the uniform
separation of every angular density copy remains the assumption for
the finite reconstruction below.  Equalities of germs, and their
translation periods, are understood in $L^1_{\mathrm{loc}}(\R^2)$
for each angle.  Thus they are independent of representatives of
the unknown density.

\begin{theorem}[Rigidity from one resolved obstacle]
\label{thm:single-resolved-obstacle}
Retain the single-law experiment, the locally finite convex
configuration, and the boundary regularity of
Section~\ref{sec:single-law-rigidity}, but replace
\eqref{eq:single-footprint-separation} by
\begin{equation}\label{eq:single-one-body-separation}
 \diam A\le\Delta<d,\qquad
       \diam C_*>\Delta\quad\hbox{for at least one }C_*\in\mathcal C.
\end{equation}
No lower width bound is required for the other obstacles, and the
identity or position of $C_*$ is not supplied.  The directional
germ determines $(\mathcal O_0,A_0,j_0)$ of
\eqref{eq:single-canonical-parameters}.  Its complete observational
fiber is \eqref{eq:single-law-fiber}; it determines the forward
collision law at every finite command length and the full obstacle
period group as in \eqref{eq:single-germ-periods}.

The footprint and density are recovered by
\eqref{eq:single-density-recovery} from any angular support component
of minimum area, where the minimum is taken over all components and
all angles.  This minimum is attained and equals $\area A$.
The occupation field is then the unique bounded continuous function
with $\inf v=0$ and distributional derivatives
\begin{equation}\label{eq:single-odd-occupation}
 \partial_{e_1}v=q_0-q_\pi,
 \qquad \partial_{e_2}v=q_{\pi/2}-q_{3\pi/2}.
\end{equation}
If $P$ ranges over the closures of the connected components of
$\{v>0\}$, the support functions of the recovered obstacles are
\begin{equation}\label{eq:single-occupation-cancellation}
                 h_{C_P^0}=h_P-h_{-A_0},
 \qquad \mathcal O_0=\bigcup_P C_P^0.
\end{equation}
\end{theorem}
\begin{proof}
The short-flight and angular identities
\eqref{eq:single-flux-boundary}--\eqref{eq:single-angular-field}
remain valid: their derivations do not use a comparison between
footprint size and obstacle width.  At each angle, different
obstacles contribute pairs of density supports whose mutual
distances are at least $d-\Delta>0$.  Within one pair the two
supports are distinct translates of $-A$; they may intersect.
Consequently each connected component of the angular measure
support is either one translate of $-A$ or the union of two
distinct translates.  Every component in the first case has area
$\area A$, while every component in the second has strictly larger
area.  Indeed, if two translates differ by $u\ne0$, their support
values in direction $u/|u|$ differ by $|u|$; the translate with
larger support contains a nonempty open cap outside the other.

Choose a diameter pair $p,q$ of $C_*$.  The maximality of
$|p-q|$ implies that $p$ and $q$ have outward normals
$n=(p-q)/|p-q|$ and $-n$, respectively.  For example,
$n\cdot(y-q)\le|y-q|\le|p-q|$ for every $y\in C_*$,
and the reverse endpoint is treated in the same way.
At the corresponding angular slice, $p-A$ and $q-A$ have
distance at least $|p-q|-\Delta>0$.  They are also separated
from every other obstacle's supports.  Both are therefore
individual support components of area $\area A$.
This proves that the global minimum is attained, and that each
minimizer is a single copy $K=c-A$.  On such a component
$S_\theta=rj(c-\,\cdot\,)$ for one $r>0$.
Its mass is $r$ and its Steiner point is $c-s(A)$, proving
\eqref{eq:single-density-recovery} exactly as before.

For completeness, define the physical occupation
\[
                 v(x)=\int_A j(z)\one_{\mathcal O}(x+z)\dd z.
\]
Since $(-a)_+-(a)_+=-a$, the flux identity gives
\begin{equation}\label{eq:single-odd-flux}
 q_\theta-q_{\theta+\pi}
   =-\sum_C\int_{\partial C}
          j(y-\,\cdot\,)\,n_C(y)\cdot n_\theta\dd s(y)
   =\partial_{n_\theta}v
\end{equation}
in distributions.  The last equality follows by testing against a
compactly supported smooth function and applying the divergence
theorem on each relevant body.  Fubini's theorem is applicable,
and only finitely many bodies meet the translated test support.
Translation continuity of $j$ in $L^1$ gives
\[
 \|v(\,\cdot+u)-v\|_\infty
                    \le\|j(\,\cdot-u)-j\|_{L^1}\longrightarrow0.
\]
Thus $v$ is bounded and continuous.  Positivity of $j$ almost
everywhere on $\operatorname{int}A$ and convexity give
\begin{equation}\label{eq:single-resolved-positivity}
        \{v>0\}=\bigcup_C\operatorname{int}(C+(-A)).
\end{equation}
Indeed, a positive integral is equivalent to a positive-area
intersection of $x+A$ with a body, and hence to an intersection of
their interiors.  The closures $C+(-A)$ have mutual gaps at least
$d-\Delta$.  A boundary point of any one of them lies in none of
these interiors and has zero occupation.  Therefore $\inf v=0$.
Two bounded continuous functions satisfying
\eqref{eq:single-odd-occupation} differ by a distribution with
zero gradient, hence by a constant.  This elementary assertion
also follows by convolution with smooth approximate identities:
every convolved difference has zero classical gradient and is
constant, and local uniform convergence gives the assertion for
the original difference.  The two zero-infimum normalizations
fix that constant.  Thus the germ determines $v$ uniquely.

The connected-component closures in
\eqref{eq:single-resolved-positivity} are exactly $P=C+(-A)$.
Writing $A=A_0+s(A)$ yields
$P=(C-s(A))+(-A_0)$.  Additivity of support functions proves
\eqref{eq:single-occupation-cancellation} and recovers every
obstacle, including those whose two angular copies overlap.
Equal germs therefore give equal canonical triples; conversely
a common translation couples every collision bit.  This proves
the claimed fiber and completion of the full forward law.

Finally, if $u$ is a period of the germ, then
$v(\,\cdot+u)$ has the same distributional gradient and the same
zero infimum as $v$.  Uniqueness gives $v(\,\cdot+u)=v$.
Translation by $u$ therefore permutes the sets $C+(-A)$;
cancelling the common footprint support proves that it permutes
the obstacles.  The converse follows from translation covariance
of the raw means.  The common period group is discrete, since a
nonzero period moves a compact component to another component at
distance at least $d$; the rank-two criterion follows as in
Corollary~\ref{cor:single-germ-periods}.
\end{proof}
